\chapter{Boolean Analysis --- Threshold Functions --- Gaussian Polar}
%%% generated by the blueprint pipeline from TCSlib.BooleanAnalysis.ThresholdFunctions.GaussianPolar
\section{Overview}
Planar Gaussian density separates into radial and angular factors. Cone masses and the angular overlap of two halfplanes give the geometric formula behind Gaussian quadrant probabilities.

\section{Declarations}

\begin{definition}[Planar standard Gaussian density]
\lean{BooleanAnalysis.ThresholdFunctions.gaussianDensity2}
\leanok
For $z=(x,y)\in\mathbb R^2$, the planar Gaussian density is $d_2(x,y)=(2\pi)^{-1}\exp(-(x^2+y^2)/2)$.
\end{definition}

\begin{lemma}[Product of scalar Gaussian densities]
\lean{BooleanAnalysis.ThresholdFunctions.gaussianPDF_mul}
\leanok
\uses{BooleanAnalysis.ThresholdFunctions.gaussianDensity2}
\difficulty{4}
For all real $x,y$, let $\varphi(t)=(2\pi)^{-1/2}e^{-t^2/2}$ and $d_2(x,y)=(2\pi)^{-1}e^{-(x^2+y^2)/2}$. Then $\varphi(x)\varphi(y)=d_2(x,y)$, with both sides regarded as nonnegative extended real values.
\end{lemma}

\begin{lemma}[Density formula for the product Gaussian law]
\lean{BooleanAnalysis.ThresholdFunctions.gaussian_prod_apply}
\leanok
\uses{BooleanAnalysis.ThresholdFunctions.gaussianDensity2, BooleanAnalysis.ThresholdFunctions.gaussianPDF_mul}
\difficulty{3}
Let $\gamma$ be the standard real Gaussian probability measure. For every measurable $A\subseteq\mathbb R^2$, $(\gamma\otimes\gamma)(A)=\int_A(2\pi)^{-1}e^{-(x^2+y^2)/2}\,dx\,dy$, where the integral is nonnegative and extended-valued.
\end{lemma}

\begin{lemma}[Normalization of the radial Gaussian factor]
\lean{BooleanAnalysis.ThresholdFunctions.lintegral_radial}
\leanok
\difficulty{5}
Let $f(r)=r\exp(-r^2/2)$ for $r>0$. Its nonnegative extended integral satisfies $\int_{(0,\infty)}f(r)\,dr=1$.
\end{lemma}

\begin{lemma}[Gaussian cone mass from angular measure]
\lean{BooleanAnalysis.ThresholdFunctions.gaussian_cone_lintegral}
\leanok
\uses{BooleanAnalysis.ThresholdFunctions.gaussianDensity2, BooleanAnalysis.ThresholdFunctions.lintegral_radial}
\difficulty{6}
Let $A\subseteq\mathbb R^2$ and $\Theta\subseteq\mathbb R$ be measurable sets such that for all real $r>0$ and $\theta$, $(r\cos\theta,r\sin\theta)\in A$ if and only if $\theta\in\Theta$. Then $\int_A(2\pi)^{-1}e^{-(x^2+y^2)/2}\,dx\,dy=(2\pi)^{-1}\operatorname{Leb}(\Theta\cap(-\pi,\pi))$, with the integral and angular Lebesgue measure in the nonnegative extended reals.
\end{lemma}

\begin{lemma}[Nonpositive cosine on a centered period]
\lean{BooleanAnalysis.ThresholdFunctions.cos_nonpos_iff_of_mem_Ioo}
\leanok
\difficulty{4}
For every real $\theta$ with $-\pi<\theta<\pi$, $\cos\theta\le0$ if and only if $\theta\le-\pi/2$ or $\pi/2\le\theta$.
\end{lemma}

\begin{lemma}[Nonpositive cosine on a negative angular interval]
\lean{BooleanAnalysis.ThresholdFunctions.cos_nonpos_iff_of_le_neg_pi_div_two}
\leanok
\difficulty{4}
For every real $u$ with $-2\pi<u\le-\pi/2$, $\cos u\le0$ if and only if $-3\pi/2\le u$.
\end{lemma}

\begin{lemma}[Nonpositive cosine with a closed left endpoint]
\lean{BooleanAnalysis.ThresholdFunctions.cos_nonpos_iff_of_neg_pi_div_two_le}
\leanok
\difficulty{4}
For every real $u$ with $-\pi/2\le u<\pi$, $\cos u\le0$ if and only if $\pi/2\le u$ or $u=-\pi/2$.
\end{lemma}

\begin{lemma}[Measurability of a two-halfplane angle set]
\lean{BooleanAnalysis.ThresholdFunctions.measurableSet_twoHalfplaneAngles}
\leanok
\difficulty{3}
For every real $\alpha$, the set $\{\theta\in\mathbb R:\cos\theta\le0\ \text{and}\ \cos(\theta-\alpha)\le0\}$ is Borel measurable.
\end{lemma}

\begin{lemma}[Angular overlap of two halfplanes]
\lean{BooleanAnalysis.ThresholdFunctions.angular_measure_two_halfplanes}
\leanok
\uses{BooleanAnalysis.ThresholdFunctions.cos_nonpos_iff_of_le_neg_pi_div_two, BooleanAnalysis.ThresholdFunctions.cos_nonpos_iff_of_mem_Ioo, BooleanAnalysis.ThresholdFunctions.cos_nonpos_iff_of_neg_pi_div_two_le, BooleanAnalysis.ThresholdFunctions.measurableSet_twoHalfplaneAngles}
\difficulty{6}
For every real $\alpha\in[0,\pi]$, the Lebesgue measure of $\{\theta\in(-\pi,\pi):\cos\theta\le0,\ \cos(\theta-\alpha)\le0\}$ equals $\pi-\alpha$.
\end{lemma}
